% =====================================================================
% paper/sections/variance_mitigation_comparison.tex
%
% Audited 2026-04-25. Earlier drafts included projected head-to-head
% results for named variance-mitigation methods. Those rows were not measured
% in this repository and several source citations could not be verified. This
% appendix now records rerun-required status instead of presenting projections.
% =====================================================================

\section{Variance-Mitigation Claims: Re-Run Required}
\label{app:variance-mitigation}

A prior version of this appendix attempted to answer reviewer questions about
variance-aware GRPO variants by reporting projected rows for AERO, CPPO,
NGRPO, Scaf-GRPO, MC-GRPO, GIFT, AReaL, and ES. Those rows were not produced by
completed training runs in this repository. They combined our baseline GRPO
trace with literature-derived or synthetic deltas and therefore cannot support
claims about held-out reward, collapse rate, time-to-collapse, or ZVF
predictiveness.

The following claims are marked rerun-required in the paper:
\begin{itemize}
  \item that any variance-mitigation method was evaluated head-to-head against
  our Qwen3-8B / GSM8K baseline at 5 seeds and 100 steps;
  \item that all surveyed methods reduce ZVF or extend time-to-collapse under
  our protocol;
  \item that Scaf-GRPO, AERO, CPPO, NGRPO, or any related method is the best
  mitigation for our traces;
  \item that ZVF remains predictive under compensation-style mitigations or
  weakens under suppression-style mitigations based on our data.
\end{itemize}

\subsection{What Remains Supported}
\label{sec:variance-survey}

The supported result is narrower: baseline GRPO runs in this artifact expose
per-step reward variance, ZVF, and gradient-utilization diagnostics. These
diagnostics can be used to design future variance-mitigation experiments, but
they do not by themselves validate any external method.

\subsection{Required Re-Run Protocol}
\label{sec:variance-integration}

Before this paper can make comparative claims about variance mitigation, each
candidate method must be implemented as executable training code and run under
the same protocol as the baseline:
\begin{enumerate}
  \item same checkpoint, tokenizer, sampler, reward function, LoRA adapter,
  optimizer, and evaluation harness;
  \item at least the same 5-seed Qwen3-8B / GSM8K sweep used for baseline
  comparisons;
  \item complete per-step logs for reward, ZVF, group size, KL/entropy where
  available, and held-out evaluation;
  \item a claim table generated directly from those logs, not from dry-run
  projections or literature deltas.
\end{enumerate}

\subsection{Retired Tables}
\label{sec:variance-head2head}
\label{sec:variance-honesty}
\label{sec:variance-zvf-stress}
\label{tab:variance-axes}
\label{tab:variance-head2head}
\label{tab:variance-zvf-rho}

The former comparison-axis table, head-to-head result table, and ZVF
predictiveness table are intentionally omitted. Their labels are retained only
so historical cross-references fail softly during transition; they must not be
interpreted as evidence.

\subsection{Summary}
\label{sec:variance-summary}

This appendix now makes one defensible statement: variance mitigation is a
future-work item. The current artifact supports measuring ZVF under baseline
GRPO, not ranking unrun mitigation methods.
